\documentclass[12pt,a4paper]{article}

\usepackage[T1]{fontenc}
\usepackage[utf8]{inputenc}
\usepackage[brazil]{babel}
\usepackage{lmodern}
\usepackage{geometry}
\usepackage{hyperref}
\usepackage{enumitem}
\usepackage{array}

\geometry{margin=2.5cm}
\hypersetup{
  colorlinks=true,
  linkcolor=blue,
  urlcolor=blue
}

\title{Plano de Aula: SQL --- Consultas e Junções}
\date{}

\begin{document}
\maketitle

\noindent
\textbf{Duração:} 2 horas (120 minutos)\\
\textbf{Público-alvo:} Iniciantes\\
\textbf{Recursos Necessários:} Laboratório com computadores, SGBD e cliente SQL, projetor (opcional), quadro e marcadores.

\section{Objetivos da Aula}
\begin{itemize}[leftmargin=*,itemsep=0.25em]
  \item Utilizar SQL para recuperação e manipulação de informações.
  \item Escrever consultas com SELECT, WHERE, ORDER BY e LIMIT (quando aplicável).
  \item Aplicar agregações (COUNT, SUM, AVG) com GROUP BY e HAVING.
  \item Realizar junções (JOIN) entre tabelas usando PK/FK.
\end{itemize}

\section{Metodologia}
Prática guiada: o professor demonstra e os alunos resolvem exercícios incrementais. O foco é transformar perguntas do mundo real em consultas e validar resultados com dados de exemplo.

\section{Cronograma e Conteúdo Detalhado (120 Minutos)}
\subsection*{Parte 1: Revisão Rápida (10 minutos)}
\begin{itemize}[leftmargin=*,itemsep=0.35em]
  \item Relembrar DDL/DML e garantir tabelas e dados de exemplo.
  \item Conectar a ideia de ``query'' com perguntas do usuário.
\end{itemize}

\subsection*{Parte 2: SELECT e Filtros (30 minutos)}
\begin{itemize}[leftmargin=*,itemsep=0.35em]
  \item SELECT * vs. selecionar colunas específicas.
  \item WHERE com comparações e AND/OR.
  \item ORDER BY para ordenar resultados.
\end{itemize}

\subsection*{Parte 3: Agregações (20 minutos)}
\begin{itemize}[leftmargin=*,itemsep=0.35em]
  \item COUNT/SUM/AVG e uso de GROUP BY.
  \item HAVING para filtrar grupos.
  \item Analogia: ``agrupando por categorias e somando totais''.
\end{itemize}

\subsection*{Parte 4: JOIN --- Cruzando Tabelas (35 minutos)}
\begin{itemize}[leftmargin=*,itemsep=0.35em]
  \item Por que JOIN existe: dados normalizados em tabelas diferentes.
  \item INNER JOIN entre CLIENTE e PET usando ID\_Cliente.
  \item Erros comuns: duplicação por relacionamento N:M e chaves erradas.
\end{itemize}

\subsection*{Parte 5: Atividade em Duplas --- Perguntas \textrightarrow\ SQL (20 minutos)}
\begin{itemize}[leftmargin=*,itemsep=0.35em]
  \item Resolver um roteiro de consultas:
  \begin{itemize}[leftmargin=*,itemsep=0.25em]
    \item ``Listar todos os pets de um cliente pelo nome.''
    \item ``Contar quantos pets cada cliente possui.''
    \item ``Listar clientes que têm ao menos 2 pets.''
    \item ``Listar pets ordenados por tipo e nome.''
  \end{itemize}
  \item Conferir resultados e discutir diferentes soluções.
\end{itemize}

\subsection*{Parte 6: Encerramento (5 minutos)}
\begin{itemize}[leftmargin=*,itemsep=0.35em]
  \item Ponte: ``Próxima aula: qualidade do esquema --- dependências funcionais e normalização.''
\end{itemize}

\end{document}

